\subsection{为什么明知道该做什么，却还是启动不了？}

短答案是：知道目标状态，并不等于具备完成转变的力量。
即使你完全理解目标状态，如果框架惯性仍然很强、当前障碍高于你可用的意志能量，或者决策窗口在任何稳定动作发生前就已经缩小，你仍然会卡住。
实践上，解决方案不是继续和自己讲道理，而是让第一步更小、更清楚、更早发生。

\begin{enumerate}[nosep]
  \item 用一句话写出目标状态，不要附加额外解释。
  \item 识别那个让旧框架持续活着的当前状态。
  \item 找出你现在能改变的最小障碍项：模糊性、准备成本、缺失工具，还是地点错配。
  \item 通过选定一个具体时间和地点，为第一步打开决策窗口。
  \item 把第一步缩减成一个可见动作，让它消耗的意志能量小于整项任务。
  \item 第一动作发生后，立刻用一个稳定动作托住新状态，避免旧框架重新长回来。
\end{enumerate}

\begin{remark}
假设你知道自己应该写报告，却始终无法开始。
有用的问题不是“我是不是太懒”，而是“哪一个第一动作能把障碍降到足以让新框架启动？”
如果答案是“打开文件、写标题、坐到桌前”，那它就是正确开端。
这套做法的价值，在于它把模糊意图改造成了一个被设计过的转变，而不是一场私人道德测验。
\end{remark}